\documentclass[10pt, a4paper]{article}
\usepackage[utf8]{inputenc}
\usepackage[T1]{fontenc}
\usepackage{amsmath, amssymb}
\usepackage{multicol}
\usepackage[margin=0.4in]{geometry}
\usepackage{booktabs}
\usepackage{enumitem}
\usepackage{titlesec}
\usepackage{fancyhdr}

% Compact list settings
\setlist{nosep}

% Compact section titles
\titleformat{\section}{\large\bfseries}{\thesection}{1em}{}
\titlespacing*{\section}{0pt}{5pt}{2pt}
\titleformat{\subsection}{\bfseries}{\thesubsection}{1em}{}
\titlespacing*{\subsection}{0pt}{3pt}{1pt}

% Custom commands for math
\newcommand{\amortized}{\hat{c}_i}
\newcommand{\actual}{c_i}
\newcommand{\potential}{\Phi}

\pagestyle{fancy}
\fancyhf{}
\lhead{Analysis of Algorithms: Amortized Analysis}
\rhead{Cheat Sheet}
\cfoot{\thepage}

\begin{document}

\begin{multicols*}{2}

% ---------------------------------------------------------
\section{1. Core Concepts}
% ---------------------------------------------------------
Amortized analysis guarantees the average performance of each operation in the worst case sequence of $n$ operations.

\subsection*{A. Aggregate Analysis}
Computes the total cost $T(n)$ for a sequence of $n$ operations directly.
\begin{itemize}
    \item \textbf{Amortized Cost:} $T(n) / n$.
    \item If total cost is $O(n)$, amortized cost is $O(1)$.
\end{itemize}

\subsection*{B. Accounting Method (The Banker's View)}
Assign an \textbf{amortized cost} ($\amortized$) to each operation.
\begin{itemize}
    \item If $\amortized > \actual$ (actual cost): Surplus stored as ``credit'' in the data structure.
    \item If $\amortized < \actual$: Use stored credit to pay for the operation.
    \item \textbf{Condition:} Total Credit $\ge 0$ at all times.
    \item \textbf{Conclusion:} $\sum_{i=1}^n \amortized \ge \sum_{i=1}^n \actual$.
\end{itemize}

\subsection*{C. Potential Method (The Physicist's View)}
Define a potential function $\potential(S)$ mapping the data structure state $S$ to a real number.
\begin{itemize}
    \item \textbf{Formula:} $d_i = \actual + \potential(S_i) - \potential(S_{i-1})$
    \item Where $d_i$ is amortized cost, $\actual$ is actual cost.
    \item \textbf{Condition:} $\potential(S_n) \ge \potential(S_0)$ (usually $\potential(S_0) = 0$ and $\potential(S_i) \ge 0$).
    \item \textbf{Conclusion:} $\sum d_i$ is an upper bound on $\sum \actual$.
\end{itemize}

% ---------------------------------------------------------
\section{2. Multipop Stack}
% ---------------------------------------------------------
\textbf{Operations:} Push, Pop, Multipop($S, k$).\\
\textbf{Naive Analysis:} Worst case per op is $O(n)$, total $O(n^2)$.\\
\textbf{Amortized Result:} Total cost $\le 2n$, Amortized $O(1)$.

\subsection*{Aggregate Proof}
Each object pushed once can be popped at most once.
\[ \text{Total Pops} \le \text{Total Pushes} \le n \]
\[ \text{Total Cost} \le n (\text{Push}) + n (\text{Pop}) = 2n \]

\subsection*{Accounting Scheme (Pricing)}
\begin{itemize}
    \item \textbf{Push:} Charge \textbf{\$2}. (\$1 for actual push, \$1 saved for future pop).
    \item \textbf{Pop / Multipop:} Charge \textbf{\$0}. (Pay using the saved \$1).
    \item \textbf{Invariant:} Bank Balance = Number of items in stack $\ge 0$.
\end{itemize}

% ---------------------------------------------------------
\section{3. Dynamic Stack (Doubling)}
% ---------------------------------------------------------
\textbf{Strategy:} Start size 1. If $S.num == S.size$, allocate new array of size $2 \cdot S.size$, copy elements, then push.

\subsection*{Complexity}
\begin{itemize}
    \item \textbf{Naive:} $O(n^2)$ (pessimistic).
    \item \textbf{Actual Total:} $< 3n$.
    \item \textbf{Amortized:} $O(1)$ per operation.
\end{itemize}

\subsection*{Aggregate Calculation}
Expensive expansion happens only when $i-1$ is a power of 2.
\[ \sum_{i=1}^n c_i = n + \sum_{j=0}^{\lg n} 2^j < 3n \]

\subsection*{Potential Function (Crucial)}
\[ \boldsymbol{\potential(S) = 2 \cdot S.num - S.size} \]
\begin{itemize}
    \item \textbf{Property:} $\potential(S_i) \ge 0$ always (since table is at least $1/2$ full).
    \item \textbf{Derivation ($d_i = 3$):}
    \begin{itemize}
        \item \textit{No Expansion:} $d_i = 1 + (2(k+1)-sz) - (2k-sz) = 3$.
        \item \textit{Expansion ($sz \to 2sz$):} $c_i = k+1$. \\
        $d_i = (k+1) + (2(k+1)-2sz) - (2k-sz) = 3$.
        \\ (Note: at overflow $k=sz$).
    \end{itemize}
\end{itemize}

% ---------------------------------------------------------
\section{4. Dynamic Stack (Contraction)}
% ---------------------------------------------------------
To support Pop operations while maintaining $O(1)$ amortized cost and $O(n)$ space.

\subsection*{Rules}
\begin{itemize}
    \item \textbf{Expansion:} Double size when Full ($num = size$).
    \item \textbf{Contraction:} Halve size when \textbf{1/4 Full} ($num = size/4$).
    \item \textbf{Note:} Do NOT shrink at 1/2 full. This causes ``thrashing'' (oscillation) leading to $O(n^2)$ complexity.
\end{itemize}
\subsection*{Result}
\begin{itemize}
    \item \textbf{Load Factor:} Always between $1/4$ and $1$.
    \item \textbf{Amortized Cost:} $O(1)$ constant time.
\end{itemize}

% ---------------------------------------------------------
\section{5. Summary Table}
% ---------------------------------------------------------
\begin{center}
\renewcommand{\arraystretch}{1.5}
\begin{tabular}{@{}p{0.2\textwidth} p{0.12\textwidth} p{0.14\textwidth}@{}}
\toprule
\textbf{Problem} & \textbf{Amortized} & \textbf{Key Formula} \\
\midrule
Multipop Stack & $O(1)$ & Push Charge $= 2$ \\
Dynamic Stack & $O(1)$ & $\Phi = 2 \cdot num - size$ \\
Binary Counter & $O(1)$ & Flip $0 \to 1$ Charge $= 2$ \\
\bottomrule
\end{tabular}
\end{center}

\section{6. Exam Tips}
\begin{itemize}
    \item \textbf{Designing $\Phi$:} Ensure $\Phi(0) = 0$ and $\Phi(i) \ge 0$.
    \item \textbf{Designing Costs:} Overcharge cheap operations (Push) to pay for expensive ones (Resize/Multipop).
    \item \textbf{Common Mistake:} Confusing \textit{Worst-Case per operation} (which can be high, e.g., $O(n)$) with \textit{Amortized Cost} (which is average $O(1)$).
\end{itemize}

\end{multicols*}
\end{document}